\documentclass[11pt]{article}
\usepackage[utf8]{inputenc}
\usepackage[T1]{fontenc}
\usepackage[italian]{babel}
\usepackage{geometry}
\geometry{a4paper, margin=1in}
\usepackage{xcolor}
\usepackage{listings}
\usepackage{hyperref}

\usepackage{tcolorbox}
\tcbuselibrary{breakable}
\begin{document}

\begin{tcolorbox}[colback=blue!5!white,colframe=blue!75!black,title=\textbf{DOMANDA}, breakable]
\texttt{form-post.html}
\begin{lstlisting}[language=HTML, basicstyle=\ttfamily\small]
<!DOCTYPE html><html>
<body>
<h2>The method Attribute</h2>
<p>This form will be submitted using the GET method:</p>
<form action="/action" target="_blank" method="post">
<label for="fname">First name:</label><br>
<input type="text" id="fname" name="fname" value="John"><br>
<label for="lname">Last name:</label><br>
<input type="text" id="lname" name="lname" value="Doe"><br><br>
<input type="submit" value="Submit">
</form>
<p>After you submit, notice that the form values is visible in the address bar
of the new browser tab.</p>
</body>
</html>
\end{lstlisting}

Esercizio
\begin{itemize}
    \item Eseguire il server web \texttt{serverHTTP.c} e con il browser preferito aprire il file \texttt{form-post.html}
    \item Cosa si vede?
    \item Provare ad analizzare il contenuto della connessione \texttt{TCP} con Wireshark.
    \item Cosa si vede?
\end{itemize}
\end{tcolorbox}

\begin{tcolorbox}[colback=green!5!white,colframe=green!75!black,title=\textbf{RISPOSTA}, breakable]
\textbf{Soluzione Esercizio 24:} Web form POST

\textbf{Operazioni Preliminari}
\begin{enumerate}
    \item Sono stati nuovamente copiati i file base del server (\texttt{serverHTTP.c}) e creata la cartella per l'esercizio.
    \item È stato creato il file "\texttt{form-post.html}" modificando l'URL di destinazione all'interno dell'attributo \texttt{action}, impostandolo a "\texttt{http://127.0.0.1:8000/action}" affinché punti al web server locale da noi lanciato. (Nota: Il testo pre-impostato all'interno del file HTML del professore conteneva un refuso, dicendo erroneamente che i valori sarebbero stati visibili nella barra degli indirizzi, ma essendo un POST questo non è vero!).
\end{enumerate}

\textbf{Risposta alla Domanda 1: "Cosa si vede (sul browser)?"}\\
Dopo aver cliccato il pulsante "Submit", il browser naviga verso la pagina di destinazione (restituendo il consueto messaggio "Hello World").\\
La differenza fondamentale rispetto all'esercizio precedente è nella BARRA DEGLI INDIRIZZI: questa volta l'URL si presenterà "pulito", mostrando esclusivamente \texttt{http://127.0.0.1:8000/action}.\\
Nessun dato (nome e cognome) viene appeso all'URL, i parametri sono completamente invisibili.

\textbf{Risposta alla Domanda 2: "Cosa si vede (analizzando il TCP con Wireshark)?"}\\
Catturando la connessione su Wireshark (o analizzando l'output che viene stampato sul terminale dal nostro server C), notiamo la struttura della richiesta \texttt{HTTP} POST:

\begin{lstlisting}[basicstyle=\ttfamily\small]
POST /action HTTP/1.1
Host: localhost:8000
User-Agent: ...
Content-Length: 20
Content-Type: application/x-www-form-urlencoded

fname=John&lname=Doe
\end{lstlisting}

A differenza del GET, qui possiamo vedere:
\begin{enumerate}
    \item Il metodo \texttt{HTTP} è cambiato in \texttt{POST}.
    \item I parametri (\texttt{fname=John\&lname=Doe}) non sono più nella prima riga dell'URI, ma si trovano isolati all'interno del "Corpo" (Body/Payload) della richiesta, separati dalle intestazioni tramite una riga vuota.
    \item Proprio a causa della presenza del Body, il client (il browser) ha dovuto specificare al server la dimensione in byte di quest'ultimo usando l'intestazione \texttt{Content-Length: 20} (che è il numero esatto di caratteri di "\texttt{fname=John\&lname=Doe}") e il tipo di codifica \texttt{Content-Type: application/x-www-form-urlencoded}.
\end{enumerate}

Guardando il codice \texttt{serverHTTP.c}, si nota proprio la presenza di una logica "\texttt{if(strcmp(method, "POST")==0)}" che intercetta i POST e utilizza il valore catturato da "Content-Length" per capire quanti caratteri ("byte") leggere per poter estrapolare dal corpo della richiesta l'intero modulo compilato dall'utente.

\textbf{Comandi Utili per Riprodurre l'Esercizio:}
\begin{lstlisting}[language=bash, basicstyle=\ttfamily\small]
# Compilare ed eseguire il server
gcc serverHTTP.c network.c -o serverHTTP
./serverHTTP

# Per testare l'invio via terminale e simulare l'HTML form con metodo POST:
curl -v -d "fname=John&lname=Doe" "http://localhost:8000/action"

# Wireshark (Filtrare per traffico HTTP sulla porta 8000 e guardare la stringa TCP):
sudo tshark -i lo -f "tcp port 8000" -z follow,tcp,ascii,0
\end{lstlisting}
\end{tcolorbox}

\end{document}
